\section*{Chapter \ref{ch:rc:counterparty-exposure} --- Counterparty Exposure}

\begin{solution}{exo:rc:counterparty-exposure:1}
A swap's value at $t$ is roughly the rate move to $t$ times the remaining annuity: the first
grows like $\sqrt t$, the second shrinks to zero, so the product peaks in between (about a
third of the way for a swap). A cross-currency swap with a final exchange keeps the full
notional's FX risk until maturity, so its exposure keeps growing with $\sqrt t$.
\end{solution}

\begin{solution}{exo:rc:counterparty-exposure:2}
With netting, $\max(10-7,0) = 3$; without, the bank loses the $+10$ and must still pay the
$7$ it owes: exposure 10.
\end{solution}

\begin{solution}{exo:rc:counterparty-exposure:3}
The required collateral is $14.3-10 = 4.3$ million; the bank holds 4.0, so the change of
0.3 million is below the MTA of 0.5: no call.
\end{solution}

\begin{solution}{exo:rc:counterparty-exposure:4}
With zero threshold the collateralised exposure is the value's move over the margin period
of risk, whose standard deviation scales with the square root of the period for a diffusion;
EE is proportional to the standard deviation. The simulation gives 1.41.
\end{solution}

\begin{solution}{exo:rc:counterparty-exposure:5}
The counterparty expects to be owed more by the bank than it owes: its exposure to the bank
is larger than the bank's to it. The bank's DVA, the value of its own default to its
creditors, is computed from the ENE and is correspondingly large.
\end{solution}

\begin{solution}{exo:rc:counterparty-exposure:6}
Specific: buying protection on a country's banks from one of those banks; a put on a
company's shares bought from the company. Right way: a commodity producer that sells its
output forward to the bank owes the bank most when prices are high, which is when the
producer is strongest.
\end{solution}

\begin{solution}{exo:rc:counterparty-exposure:7}
USD~1.83 million against 1.30 million, a ratio of 1.41, close to $\sqrt2 = 1.414$.
\end{solution}

\begin{solution}{exo:rc:counterparty-exposure:8}
Margin lags: on default the bank holds the collateral of the last call met and bears the
move over the margin period of risk (plus the MTA, disputes and settlement flows); the
collateralised PFE of the chapter's \omterm{def:rc:counterparty-exposure:netting}{netting set} is still USD~6.45 million. The counterparty
needs a limit on that residual exposure, and on the risk that it stops posting.
\end{solution}

\begin{solution}{pb:rc:counterparty-exposure:1}
\textbf{1.} 3.48\%; 1.1166 dollars per euro.
\textbf{2.} USD~4.42 million, just before the forward settles at one year.
\textbf{3.} USD~626\,224 over one year and 545\,826 over five.
\textbf{4.} The forward: EUR~20 million at 8\% volatility is larger risk than the swap in its
first year, when little of the rate uncertainty has accumulated.
\textbf{5.} Posting cash daily needs liquidity and operations that most corporates do not
have, and hedge accounting prefers stable cash flows.
\textbf{6.} Peak PFE USD~3.63 million; EPE USD~199\,002 over one year and 97\,606 over five.
\textbf{7.} The forward settles at one year: its value leaves the \omterm{def:rc:counterparty-exposure:netting}{netting set} as a payment
while the collateral held still reflects it for one margin period of risk; if the bank had
posted collateral against it, that collateral is exposed until returned.
\textbf{8.} Peak PFE USD~4.38 million; life EPE USD~544\,793.
\textbf{9.} The \omterm{def:rc:counterparty-exposure:netting}{netting set}'s value rarely exceeds USD~5 million, so the threshold is seldom
reached and almost no collateral would move.
\textbf{10.} Cash or securities to post, a treasury able to meet daily calls, and possibly a
credit line for margin.
\textbf{11.} No: its peak PFE of USD~4.42 million exceeds the limit.
\textbf{12.} New trades that add exposure at the one-year peak need limit headroom; trades that
reduce it (the client selling euros forward) are welcome.
\textbf{13.} A CSA, a shorter forward, settling the forward in two parts, or a break clause;
payment-versus-payment settlement removes the delivery risk on the day.
\textbf{14.} The bank's short euro forward loses about $20\text{ million}\times0.11 = \text{USD}~2.2$
million: its exposure to the client falls, the client's exposure to the bank rises by as much.
\textbf{15.} If the euro falls the forward (bank sells euros at 1.1166) gains for the bank: the
bank is exposed to the client.
\textbf{16.} Corporates generate the hedging business and have few alternatives; the bank
charges for the exposure (chapters~18 to~20) and limits it.
\textbf{17.} Only the collateralised figures change, by about $\sqrt2$ on the residual exposure;
the uncollateralised profile is unchanged.
\textbf{18.} Protection up to the threshold: the exposure below it stays uncollateralised and
must be priced and limited.
\textbf{19.} Named result: \emph{the uncollateralised corporate}: peak PFE USD~4.42 million and
life EPE USD~545\,826 without a CSA; USD~3.63 million and 97\,606 with the zero-threshold CSA.
\textbf{20.} Into the risk of the value moving over a few days, plus the settlement flows.
\end{solution}

\begin{solution}{iq:rc:counterparty-exposure:1}
EE: the mean positive exposure at a date, used to price CVA. PFE: a high quantile of the
exposure at a date, used for credit limits. EPE: the time average of EE, a single number
used in regulatory exposure measures and in capital.
\iqlookfor{the three definitions and their uses.}
\end{solution}

\begin{solution}{iq:rc:counterparty-exposure:2}
Common scenarios for all risk factors on a grid; revaluation of every trade on every path
and date (closed forms, grids, regression for callables), organised by trade type and
vectorised; netting and collateral per \omterm{def:rc:counterparty-exposure:netting}{netting set}; profiles stored per \omterm{def:rc:counterparty-exposure:netting}{netting set} for
the adjustments. The time goes into revaluation; the aggregation is memory-bound; both
parallelise over paths.
\iqlookfor{scenario, revaluation and aggregation, and where the cost is.}
\end{solution}

\begin{solution}{iq:rc:counterparty-exposure:3}
The time between the last collateral received and the close-out and re-hedging of the
defaulted \omterm{def:rc:counterparty-exposure:netting}{netting set}: noticing the missed call, disputes, grace periods, termination and
replacement. Ten business days is the regulatory minimum for daily-margined bilateral
trades, more for large or disputed \omterm{def:rc:counterparty-exposure:netting}{netting sets}.
\iqlookfor{the steps in the period and the regulatory floors.}
\end{solution}

\begin{solution}{iq:rc:counterparty-exposure:4}
A bank buying protection on a sovereign from that sovereign's banks; an FX forward in which
a local company buys dollars and would default after a devaluation. Measure it by linking the
counterparty's hazard to the market factors in the simulation and computing exposure
conditional on default, or by stress tests of joint scenarios.
\iqlookfor{an example and a quantitative method.}
\end{solution}

\begin{solution}{iq:rc:counterparty-exposure:5}
Its value at each simulation date depends on the future exercise decision: use regression
(Longstaff--Schwartz) on the simulated state to estimate the continuation value and the
exercise indicator on each path, then its value on the path; after exercise, the trade
becomes the underlying swap on that path.
\iqlookfor{regression revaluation and path-dependent exercise.}
\end{solution}

\begin{solution}{iq:rc:counterparty-exposure:6}
They terminated their trades under the master agreements, netted, kept collateral, and
replaced hedges in a stressed market all at once; most of the roughly 930\,000 transactions
were terminated within two months, some were still open in January 2009, and disputes over
close-out amounts lasted years, with claims above 50 billion dollars.
\iqlookfor{close-out, replacement cost and disputes.}
\end{solution}
